Прежде чем переходить к описанию экспериментов в симулируемых средах, необходимо формализовать алгоритмический жизненный цикл нейроэволюционного подхода. В данном разделе подробно рассматривается внутренняя архитектура алгоритма NEAT и механизмы его работы, которые применялись при обучении когнитивных агентов.

\subsection{Алгоритмический базис: Жизненный цикл NEAT}

В отличие от классического обучения с подкреплением, где оптимизируются только веса заранее заданной графовой структуры, NEAT управляет популяцией агентов, одновременно изменяя как веса, так и саму топологию графа. Процесс обучения представляет собой циклический эволюционный алгоритм, состоящий из нескольких ключевых этапов.

\subsubsection{Кодирование нейросети: генотип и фенотип}

Основополагающим принципом NEAT является разделение графа нейронной сети на генотип (информационный код, подверженный мутациям) и фенотип (рабочая нейронная сеть, управляющая когнитивным агентом в среде).

Генотип особи в NEAT состоит из двух списков генов:
\begin{itemize}
  \item Гены узлов: определяют входы, выходы и скрытые нейроны агента. Каждый узел может иметь свою индивидуальную функцию активации (например, ReLU, Sigmoid или Tanh), что позволяет алгоритму подбирать оптимальную нелинейность для конкретной задачи.
  \item Гены связей: описывают направленные синапсы между узлами. Каждый ген связи содержит информацию о начальном и конечном узле, текущем весе, статусе активности и, что наиболее важно, номер инновации.
\end{itemize}

Компиляция генотипа в фенотип происходит каждый раз, когда агента нужно протестировать в симуляции. Генотип подвергается мутациям и скрещиванию, а затем компилируется в конкретную архитектуру нейросети, которая управляет поведением агента в среде. Этот фенотип может быть как простой однослойной сетью, так и сложной рекуррентной архитектурой с несколькими скрытыми слоями и различными функциями активации.

\subsubsection{Отказ от слоистой архитектуры в пользу произвольных графов}

Важным концептуальным отличием алгоритма NEAT от традиционных методов глубокого обучения является полный отказ от идеи "слоев". В классических искусственных нейронных сетях, таких как многослойный перцептрон, архитектура строится строго послойно: сигнал синхронно передается от входного слоя к первому скрытому, затем ко второму и так далее, вплоть до выходного слоя. Вычисления при этом производятся структурными блоками через операции умножения матриц.

Такой отказ наделяет когнитивного агента новыми возможностями. Во-первых, один сенсор может быть соединен с моторным выходом напрямую, тогда как сигнал от соседнего датчика будет обрабатываться через сложную цепочку промежуточных узлов. Во-вторых, теперь глубина сети определяется не количеством заранее заданных слоев-матриц, а длиной максимального пути в графе от входа к выходу. Также, масштабирование графа происходит точечно. Когда алгоритм принимает решение об усложнении архитектуры, он не генерирует целый скрытый слой с массой избыточных нейронов, а лишь аккуратно расщепляет одну конкретную связь, добавляя единичный вычислительный узел именно в ту часть графа, где остро требуется дополнительная нелинейность.


\subsubsection{Скрещивание и решение проблемы «соперничающих соглашений»}


Скрещивание является одним из наиболее сложных операторов в нейроэволюции. Основная трудность классических подходов заключалась в так называемой проблеме соперничающих соглашений, проиллюстрованная на рисунке \ref{fig:competing_conventions}. Она возникает, когда две нейронные сети кодируют одну и ту же функцию разными способами (например, разным расположением скрытых нейронов). При обычном случайном скрещивании таких сетей информация в геноме «перемешивается» деструктивно, что приводит к потере функциональных блоков и резкому падению приспособленности потомков.

\begin{figure}[h]
  \centering
  \begin{tikzpicture}[
      node distance=1.5cm,
      neuron/.style={circle, draw, minimum size=6mm, thick, inner sep=0pt},
      input/.style={neuron, fill=white},
      hiddenA/.style={neuron, fill=white},
      hiddenB/.style={neuron, fill=gray!40},
      hiddenC/.style={neuron, fill=black, text=white},
      output/.style={neuron, fill=white},
      edge/.style={->, >=stealth, semithick, draw=gray!80}
  ]
  
  % --- ЛЕВАЯ НЕЙРОСЕТЬ ---
  % Входы
  \node[input, label=below:1] (I1L) at (0,0) {};
  \node[input, label=below:2] (I2L) at (1.5,0) {};
  
  % Скрытый слой
  \node[hiddenA, label=left:A] (HAL) at (-0.5, 1.5) {};
  \node[hiddenB, label=above left:B] (HBL) at (0.75, 1.5) {};
  \node[hiddenC, label=right:C] (HCL) at (2.0, 1.5) {};
  
  % Выход
  \node[output, label=above:3] (O1L) at (0.75, 3) {};
  
  % Связи (Левая сеть)
  \foreach \i in {I1L, I2L} {
      \foreach \h in {HAL, HBL, HCL} {
          \draw[edge] (\i) -- (\h);
      }
  }
  \foreach \h in {HAL, HBL, HCL} {
      \draw[edge] (\h) -- (O1L);
  }
  
  % --- ПРАВАЯ НЕЙРОСЕТЬ ---
  % Входы
  \node[input, label=below:1] (I1R) at (5,0) {};
  \node[input, label=below:2] (I2R) at (6.5,0) {};
  
  % Скрытый слой (Порядок узлов изменен)
  \node[hiddenC, label=left:C] (HCR) at (4.5, 1.5) {};
  \node[hiddenB, label=above right:B] (HBR) at (5.75, 1.5) {};
  \node[hiddenA, label=right:A] (HAR) at (7.0, 1.5) {};
  
  % Выход
  \node[output, label=above:3] (O1R) at (5.75, 3) {};
  
  % Связи (Правая сеть)
  \foreach \i in {I1R, I2R} {
      \foreach \h in {HCR, HBR, HAR} {
          \draw[edge] (\i) -- (\h);
      }
  }
  \foreach \h in {HCR, HBR, HAR} {
      \draw[edge] (\h) -- (O1R);
  }
  
  % --- ТЕКСТ ПОД РИСУНКОМ ---
  \node[align=center] at (3.25, -1.5) {
      $[A, B, C]$ \\
      $\times \ [C, B, A]$ \\
      \rule{4cm}{0.4pt} \\
      Потомки: $[A, B, A] \quad [C, B, C]$ \\
      \footnotesize (разрушение структуры: потеря информации)
  };
  
  \end{tikzpicture}
  \caption{Проблема конкурирующих соглашений}
  \label{fig:competing_conventions}
  \end{figure}

Алгоритм NEAT решает эту проблему с помощью исторических меток. Каждому новому гену (связи или узлу), появившемуся в результате структурной мутации, присваивается уникальный идентификатор, который сохраняется за этой инновацией во всей популяции. Это позволяет алгоритму «выравнивать» геномы двух родителей, точно определяя, какие части их «цифрового мозга» имеют общее происхождение, а какие являются уникальными разработками конкретной эволюционной ветки.

\subsubsection{Классификация генов при выравнивании}

При создании потомка от двух родителей алгоритм сравнивает их списки генов связей. Гены распределяются по трем категориям на основе их номеров инноваций:

\begin{enumerate}
  \item Совпадающие гены: Гены с одинаковыми номерами инноваций у обоих родителей. Это означает, что данные связи были унаследованы от общего предка. Потомок получает этот ген от одного из родителей случайным образом, а веса связей усредняются или также выбираются случайно.

  \item Разъединенные гены: Гены, которые присутствуют только у одного из родителей, но их номер инновации находится внутри диапазона номеров другого родителя. Это инновации, которые произошли в прошлом одной ветви, но отсутствуют в другой из-за того, что та ветвь пошла по иному пути развития.

  \item Избыточные гены: Гены, которые присутствуют только у одного из родителей и имеют номера инноваций, превышающие максимальный номер другого родителя. Это самые свежие структурные изменения, появившиеся совсем недавно в процессе эволюции одной из ветвей.

\end{enumerate}

Допустим, мы скрещиваем Родителя А, содержащего гены с маркерами 1, 2, 3, 8, 10 и Родителя Б с генами 1, 2, 5, 6. Алгоритм моментально определяет, что гены 1 и 2 являются совпадающими — они достались агентам от общего предка. Гены 3, 5 и 6 классифицируются как разъединенные, поскольку их номера не превышают максимальные индексы противоположных родителей (8 и 6 соответственно), но пары для них отсутствуют. В свою очередь, гены 8 и 10 у Родителя А однозначно распознаются как избыточные, так как они выходят за верхнюю границу инноваций Родителя Б.

\subsubsection{Вычисление функции приспособленности}

В контексте моделирования когнитивных агентов функция приспособленности напрямую эквивалентна кумулятивной награде $J(\theta)$ из математической формализации RL (формула \ref{eq:rl_fitness}).

На каждом поколении для каждой особи в популяции выполняется следующий алгоритм:
\begin{enumerate}

  \item Из генотипа компилируется фенотип (нейронная сеть прямого распространения или рекуррентная сеть).

  \item Агент, управляемый этой сетью, помещается в среду (например, CartPole или LunarLander из фреймворка \cite{towers2025gymnasiumstandardinterfacereinforcement}).
  
  \item Агент взаимодействует со средой до наступления терминального состояния (успех, провал или достижение лимита шагов).
  
  \item Суммарная накопленная награда, полученная от среды, присваивается особи в качестве её показателя приспособленности.
  
\end{enumerate}

Именно этот показатель определяет вероятность особи передать свои гены следующему поколению.

Для того чтобы предотвратить накопление «мусорных» структур от неудачных особей, в NEAT применяется строгое правило: разъединенные и избыточные гены наследуются только от более приспособленного родителя, имеющего более высокий показатель приспособленности. Если же родители имеют равную приспособленность, то уникальные гены передаются потомку от обоих родителей с некоторой вероятностью.

\subsubsection{Операторы эволюции: мутация и скрещивание}

За внесение новых генетических признаков в популяцию отвечают мутации. Чтобы алгоритм мог независимо развивать как параметры, так и саму структуру сети, этот механизм разделен на два принципиально разных класса.

Первый класс — параметрические мутации. Они не меняют архитектуру, а работают с уже существующим графом. С заданной вероятностью алгоритм либо слегка сдвигает веса синапсов, добавляя случайное значение из нормального распределения, либо полностью заменяет их на новые числа. Сюда же относится случайная смена функции активации у отдельного узла.

Второй класс — структурные мутации, являющиеся главным двигателем усложнения топологии. Здесь работают два оператора. Оператор добавления связи просто выбирает два случайных узла, между которыми еще нет прямого контакта, и прокладывает новый синапс. Это открывает новые маршруты для сигнала. Алгоритм берет активную связь, отключает её и ставит на её место новый скрытый нейрон. Входящая к нему связь получает вес 1.0, а исходящая наследует вес старой отключенной связи. Благодаря этому в первый момент после мутации выход сети не меняется, но у алгоритма появляется новый свободный параметр для оптимизации.

\paragraph{Скрещивание геномов} позволяет комбинировать удачные структурные блоки от разных особей. При создании потомка от двух родителей алгоритм должен корректно объединить сети, которые могут иметь совершенно разную топологию.

Здесь важную роль играют исторические маркеры. Алгоритм выстраивает гены родителей в ряд и производит выравнивание:

\begin{enumerate}

  \item Совпадающие гены, имеющие одинаковые номера инноваций у обоих родителей, наследуются потомком случайным образом от отца или от матери.

  \item Уникальные гены, присутствующие только у одного из родителей, не смешиваются вслепую. Они передаются потомку только в том случае, если принадлежат более успешному родителю (с более высоким показателем приспособленности). Если оба родителя одинаково успешны, уникальные гены могут быть унаследованы от обоих.

\end{enumerate}

Такой подход гарантирует, что потомки будут накапливать полезные топологические инновации, не разрушая при этом уже сформированные функциональные блоки своих предков.


\subsubsection{Видообразование на практике}

Новые структурные элементы часто кратковременно снижают приспособленность особи, из-за чего в стандартных эволюционных алгоритмах инновации быстро вымирают. Для защиты структурных инноваций NEAT разбивает популяцию на виды.

Разделение происходит на основе метрики генетической дистанции $\delta$, которая вычисляется по формуле:
\begin{equation}
  \label{eq:genetic_distance}
  \delta = \frac{c_1 E}{N} + \frac{c_2 D}{N} + c_3 \cdot \overline{W}
\end{equation}

где $E$ — количество excess-генов, $D$ — количество disjoint-генов, $\overline{W}$ — среднее различие в весах совпадающих генов, $N$ — нормализующий фактор (обычно принимается равным размеру большей из двух родительских сетей), а $c_1$, $c_2$, $c_3$ — коэффициенты, регулирующие важность каждого компонента.

Если дистанция $\delta$ между особью и представителем вида меньше заданного порога $\delta_{threshold}$, то особь зачисляется в этот вид. Далее применяется механизм разделения приспособленности: скорректированная приспособленность особи рассчитывается путем деления её реальной приспособленности на размер её вида.

Благодаря этому когнитивные агенты конкурируют в основном внутри своих видов. Это не позволяет одной простой, но локально успешной топологии мгновенно захватить всю популяцию, давая сложным архитектурам время на оптимизацию своих новых весов. Инкрементальный рост сложности от минимальных сетей к более сложным является ключом к эффективности NEAT при поиске оптимального мозга агента.